% p005_d 的电脑重画（③ 饱和光电流测定电路：左斜板接 A 表与电池，右竖板接滑动变阻器滑片；hν、e^-、右上角“原因”与弧线）；坐标取原图像素，1px=0.155mm
\begin{tikzpicture}[x=0.155mm, y=-0.155mm, line width=0.7pt, >={Stealth[length=2.2mm]}]
  % 左侧导线、电流表
  \draw (96,83) -- (35,85) -- (35,138);
  \draw (38,183) -- (38,262);
  \draw (39,160) ellipse (20 and 23);
  \node at (39,161) {A};
  % 左斜板
  \draw (75,52) -- (133,128);
  % 滑动变阻器及其左右导线
  \draw (38,200) -- (141,201);
  \draw (143,186) rectangle (205,205);
  \draw (205,197) -- (250,195) -- (299,193) -- (303,230) -- (303,268);
  % 电池
  \draw (36,262) -- (160,266);
  \draw (160,244) -- (160,281);
  \draw (175,238) -- (175,287);
  \draw (176,265) -- (250,271) -- (303,270);
  % 右竖板及其导线、滑片
  \draw (235,38) -- (234,118);
  \draw (235,51) -- (293,50) -- (297,70) -- (298,110) -- (296,140) -- (293,152) -- (270,157) -- (240,158) -- (215,160) -- (195,162) -- (178,158) -- (163,157);
  \draw[->] (163,157) -- (163,183);
  % 电子、光
  \node at (138,72) {$e^-$};
  \draw[->] (183,76) -- (214,74);
  \draw[->] (205,3) -- (175,34);
  \node at (240,28) {$h\nu$};
  % 右上角“原因”与弧线
  \node[anchor=west] at (284,45) {原因};
  \draw plot[smooth] coordinates {(262,0) (272,10) (285,17) (300,20) (318,19) (341,12)};
\end{tikzpicture}
